% ======================================================================
%  PART 6 -- Fourier Theory, Spectral Density and the Periodogram
% ======================================================================
\section{Fourier Theory, Spectral Density and the Periodogram}
\label{sec:spectral}

The frequency domain is the second face of a stationary process. Every
second-order property carried by the autocovariance sequence (ACVS)
$\{\acvs_\tau\}$ in the time domain is mirrored, via the Fourier transform, by
the \textbf{spectral density function} (SDF) $\sdf(f)$ in the frequency domain.
This part builds that bridge in three stages. First, pure Fourier theory:
continuous-time, discrete-time and finite-data transforms, the convolution
theorem, aliasing and the Nyquist frequency. Second, the spectral
representation of a stationary process: $\sdf(f)=\sum_\tau\acvs_\tau
e^{-2\pi i f\tau}$, Herglotz's theorem, the integrated spectrum and its
Lebesgue decomposition, and the spectral representation theorem with its
orthogonal-increment process $\{Z(f)\}$. Third, estimation: the periodogram
$\hat\sdf^{(p)}(f)=\abs{J(f)}^2$, why it is (asymptotically) unbiased but
\emph{inconsistent}, and how tapering cures bias while multitapering cures
variance. Throughout we take the sampling interval $\Delta=1$ once we reach
the spectral chapters, and we always work with ordinary frequency $f$ (not
angular frequency $\omega=2\pi f$).

% ======================================================================
\subsection{Definitions}
% ======================================================================

\begin{definition}{Continuous-time Fourier transform and its inverse}{p6-ctft}
For $g:\Reals\to\Comp$ with $g\in L^1$, the \textbf{Fourier transform}
$G$ (upper case denotes the transform; some authors write $\hat g$) is
\[
G(f)=\int_{-\infty}^{\infty} g(t)\,e^{-2\pi i f t}\,dt,\qquad f\in\Reals .
\]
If in addition $G\in L^1$, the \textbf{inverse transform} recovers $g$ for
almost all $t$:
\[
g(t)=\int_{-\infty}^{\infty} G(f)\,e^{2\pi i f t}\,df .
\]
We write $g\longleftrightarrow G$ and call $(g,G)$ a \textbf{Fourier pair}.
\end{definition}

\begin{definition}{Properties of the continuous-time transform}{p6-ctprops}
For $g,h\in L^1$ with transforms $G,H$ and scalar $\alpha$, the basic pairs are
\[
\begin{aligned}
\conj{g(t)} &\longleftrightarrow \conj{G(-f)}, &\qquad
g(\alpha t)\;(\alpha\ne0) &\longleftrightarrow \tfrac{1}{\abs\alpha}G(f/\alpha),\\
g(t+\tau) &\longleftrightarrow G(f)\,e^{2\pi i \tau f}, &\qquad
\alpha g(t)+h(t) &\longleftrightarrow \alpha G(f)+H(f).
\end{aligned}
\]
(Conjugation reflects, scaling in time inversely scales frequency, a time
shift multiplies by a phase, and the transform is linear.)
\end{definition}

\begin{definition}{Continuous-time convolution}{p6-ctconv}
For $g,h\in L^1$, the \textbf{convolution} $u=g*h$ is
\[
u(t)=\int_{-\infty}^{\infty} g(t-u)\,h(u)\,du,\qquad t\in\Reals .
\]
Convolution at a point is a weighted average of one function by the other,
centred at that point.
\end{definition}

\begin{definition}{Discrete-time Fourier transform and its inverse}{p6-dtft}
For a sequence $\{g_t\}_{t\in\Delta\Ints}$ in $\ell^1$, the
\textbf{discrete-time Fourier transform} is
\[
G(f)=\Delta\sum_{t\in\Delta\Ints} g_t\,e^{-2\pi i t f},\qquad f\in\Reals,
\]
with inverse (a Fourier series with the roles of domains swapped)
\[
g_t=\int_{-1/(2\Delta)}^{1/(2\Delta)} G(f)\,e^{2\pi i t f}\,df,\qquad t\in\Delta\Ints .
\]
A defining feature is \textbf{periodicity}: $G(f)=G(f+k/\Delta)$ for every
$k\in\Ints$, so $G$ is determined by its values on one period of length
$1/\Delta$.
\end{definition}

\begin{definition}{Discrete-time and periodic convolution}{p6-dtconv}
For $g,h\in\ell^1$ the \textbf{discrete convolution} $u=g*h$ is
\[
u_t=\Delta\sum_{u\in\Delta\Ints} g_{t-u}\,h_u,\qquad t\in\Delta\Ints,
\]
(with $\Delta=1$ the prefactor disappears). The dual object, the
\textbf{periodic convolution} $U=G*H$ of two periodic spectra, is
\[
U(f)=\int_{-1/(2\Delta)}^{1/(2\Delta)} G(f-f')\,H(f')\,df' .
\]
\end{definition}

\begin{definition}{Finite-data Fourier transform and the Dirichlet kernel}{p6-fdft}
Given observations $g_0,\dots,g_{(n-1)\Delta}$ at times
$T=\{0,\Delta,\dots,(n-1)\Delta\}$, the \textbf{finite Fourier transform} is
\[
G_n(f)=\sum_{t\in T} g_t\,e^{-2\pi i t f},\qquad f\in\Reals .
\]
It equals the full transform of $g_t$ multiplied by the indicator
$d_t=\tfrac1\Delta\ind_T(t)$, whose own transform is the \textbf{Dirichlet
kernel}
\[
\Dir_T(f)=\sum_{t\in T} e^{-2\pi i f t}
=\frac{\sin(n\pi f\Delta)}{\sin(\pi f\Delta)}\,e^{-\pi i (n-1) f\Delta}.
\]
Thus $G_n=\Dir_T * G$: observing a finite window \emph{blurs} the true
transform.
\end{definition}

\begin{definition}{Fourier frequencies}{p6-fourierfreq}
The FFT evaluates $G_n$ at the \textbf{Fourier frequencies}
\[
f_k=\frac{k}{\Delta n},\qquad k\in\Ints,\quad -\frac{n}{2}\le k\le \frac{n-1}{2}.
\]
With $\Delta=1$ these are $f_k=k/n$. They are equally spaced over one period
and are the frequencies at which distinct periodogram ordinates are
(approximately) uncorrelated.
\end{definition}

\begin{definition}{Spectral density function (discrete and continuous)}{p6-sdf}
For a discrete-time stationary process $\{X_t\}_{t\in\Ints}$ with ACVS
$\{\acvs_\tau\}\in\ell^1$, the \textbf{spectral density function} is the
discrete-time Fourier transform of the ACVS (with $\Delta=1$):
\[
\sdf(f)=\sum_{\tau\in\Ints}\acvs_\tau\,e^{-2\pi i f\tau},\qquad f\in\Reals .
\]
For a continuous-time process with $\acvs\in L^1$,
$\sdf(f)=\int_{-\infty}^{\infty}\acvs(\tau)e^{-2\pi i f\tau}\,d\tau$. By Fourier
inversion the ACVS is recovered as
$\acvs_\tau=\int_{-1/2}^{1/2}\sdf(f)e^{2\pi i f\tau}\,df$; in particular
$\acvs_0=\Var(X_t)=\int_{-1/2}^{1/2}\sdf(f)\,df$, so the SDF distributes the
total ``power'' $\Var(X_t)$ across frequencies. A valid SDF is real, even
($\sdf(-f)=\sdf(f)$, since $\acvs_{-\tau}=\acvs_\tau$), non-negative, and
periodic with period $1$.
\end{definition}

\begin{definition}{Harmonic (sinusoidal) process}{p6-harmonic}
Let $A>0$ and $\Theta\sim\mathrm{Unif}(-\pi,\pi)$ be independent. The
\textbf{harmonic process} is
\[
X_t=A\cos(\nu t+\Theta),\qquad t\in\Ints,
\]
with $\nu\in\Reals$ the oscillation frequency. Then $\E[X_t]=0$ and
\[
\Cov(X_{t+\tau},X_t)=\tfrac12\,\E[A^2]\cos(\nu\tau),
\]
so it is stationary but its ACVS does \emph{not} decay; in particular
$\{\acvs_\tau\}\notin\ell^1$ and no ordinary SDF exists. Its spectrum is a
pair of lines (jumps) at $\pm\nu/(2\pi)$.
\end{definition}

\begin{definition}{Integrated spectrum (spectral distribution function)}{p6-integspec}
By Herglotz's theorem, every ACVS admits the representation
\[
\acvs_\tau=\int_{-1/2}^{1/2} e^{2\pi i \tau f}\,dS^{(I)}(f),
\]
where $S^{(I)}$, the \textbf{integrated spectrum}, is right-continuous,
bounded, non-decreasing, with $S^{(I)}(-\tfrac12)=0$ and
$S^{(I)}(\tfrac12)=\sigma^2=\acvs_0$. It behaves like a (scaled) cumulative
distribution function over frequency. When $\{\acvs_\tau\}\in\ell^1$ there is a
density $\sdf$ with $S^{(I)}(f)=\int_{-1/2}^{f}\sdf(\lambda)\,d\lambda$, i.e.\
$\sdf=dS^{(I)}/df$, recovering Definition~\ref{def:p6-sdf}.
\end{definition}

\begin{definition}{Periodogram}{p6-periodogram}
For data $X_1,\dots,X_n$ ($T=\{1,\dots,n\}$, $\Delta=1$), the
\textbf{periodogram} is the plug-in estimator obtained by Fourier transforming
the biased sample ACVS:
\[
\hat\sdf^{(p)}(f)=\sum_{\tau\in\Ints}\hat\acvs_\tau\,e^{-2\pi i f\tau},
\qquad
\hat\acvs_\tau=\frac1n\sum_{t=1}^{n-\abs\tau}(X_t-\bar X)(X_{t+\abs\tau}-\bar X)
\]
for $\abs\tau\le n-1$ ($0$ otherwise). Equivalently (Proposition
\ref{prop:p6-periogramJ}) $\hat\sdf^{(p)}(f)=\abs{J(f)}^2$ with
$J(f)=\sqrt{1/n}\sum_{t\in T}(X_t-\bar X)e^{-2\pi i t f}$.
\end{definition}

\begin{definition}{Data taper and tapered periodogram}{p6-taper}
A \textbf{data taper} is a sequence $\{h_t\}_{t\in T}$ with
$\norm{h}_2^2=\sum_t h_t^2=1$. The \textbf{tapered Fourier transform} and
\textbf{tapered periodogram} are
\[
J_h(f)=\sum_{t\in T} h_t\,(X_t-\bar X)\,e^{-2\pi i t f},
\qquad
\hat\sdf^{(p)}_h(f)=\abs{J_h(f)}^2 .
\]
Choosing $h_t=\sqrt{1/n}$ (the constant taper) recovers the ordinary
periodogram. The normalisation $\norm{h}_2^2=1$ keeps the estimator unbiased
for white noise (Exercise~\ref{exo:p6-tunbias}).
\end{definition}

\begin{definition}{Multitaper spectral estimator}{p6-multitaper}
Given $K$ \textbf{orthonormal} tapers $h_k=\{h_{t,k}\}_{t\in T}$ with
$\sum_{t\in T} h_{t,k}h_{t,k'}=\delta_{k,k'}$, the \textbf{multitaper
estimator} averages the individual tapered periodograms,
\[
\hat\sdf^{(mt)}(f)=\frac1K\sum_{k=1}^{K}\hat\sdf^{(p)}_{h_k}(f).
\]
Its expectation is a convolution with the \textbf{average kernel}
$\overline H(f)=\tfrac1K\sum_{k=1}^K\abs{H_k(f)}^2$, and its variance is
$\approx\sdf(f)^2/K$.
\end{definition}

% ======================================================================
\subsection{Theorems \& Propositions}
% ======================================================================

\begin{theorem}{Fourier inversion (continuous and discrete time)}{p6-inversion}
\emph{(Continuous time.)} If $g,G\in L^1$ then the inverse transform recovers
$g$ for almost all $t$:
$g(t)=\int_{-\infty}^{\infty} G(f)\,e^{2\pi i f t}\,df$.
\emph{(Discrete time.)} If $\{g_t\}\in\ell^1$ then $G$ is recovered on one
period as a Fourier series with the domains swapped:
\[
g_t=\int_{-1/(2\Delta)}^{1/(2\Delta)} G(f)\,e^{2\pi i t f}\,df,\qquad
t\in\Delta\Ints .
\]
\textit{Proof idea:} substitute the forward transform into the right-hand side
and use the Fourier orthogonality relation
($\int e^{2\pi i f(t-s)}\,df$ acts as a delta), the integrability hypotheses
justifying the interchange of integration/summation.\\
\textit{Why:} inversion is what makes $(g,G)$ a genuine \emph{Fourier pair} ---
no information is lost. It is the engine behind recovering the ACVS from the SDF
($\acvs_\tau=\int_{-1/2}^{1/2}\sdf(f)e^{2\pi i f\tau}\,df$) and behind every
``time $\leftrightarrow$ frequency'' duality used below.
\end{theorem}

\begin{theorem}{Convolution theorem (continuous time)}{p6-convthm}
For $g,h\in L^1$, if $u=g*h$ then $U=G\cdot H$: the transform of a convolution
is the product of the transforms. If also $G,H\in L^1$ and $v=g\cdot h$, then
dually $V=G*H$.\\[2pt]
\textit{Proof idea:} substitute the convolution integral into the definition of
$U$, swap the order of integration (Fubini, justified by
$\iint\abs{g(x-y)h(y)}\,dx\,dy=\norm g_1\norm h_1<\infty$), and recognise the
shift factor $e^{-2\pi i s f}$ as producing $G(f)$.\\
\textit{Why:} the central computational identity of Fourier analysis ---
``multiplication in one domain is convolution in the other''. It underlies
filtering, the SDF of a filtered process, and the periodogram bias formula.
\end{theorem}

\begin{theorem}{Convolution theorem (discrete time)}{p6-dconvthm}
For $g,h\in\ell^1$ both pairs hold:
\[
h\cdot g \longleftrightarrow H*G,\qquad
h*g \longleftrightarrow H\cdot G .
\]
\textit{Proof idea:} for $u=h*g$, expand $U(f)=\sum_t u_t e^{-2\pi i t f}$,
insert $e^{-2\pi i (t-s)f}e^{-2\pi i s f}$, factor the double sum into
$\big(\sum_s h_s e^{-2\pi i s f}\big)\big(\sum_x g_x e^{-2\pi i x f}\big)$;
interchange is legal because $g,h\in\ell^1$. For $v=h\cdot g$, replace one
factor by its inverse transform and integrate.\\
\textit{Why:} the discrete analogue used to derive the finite-data transform
($G_n=\Dir_T*G$) and the periodogram expectation (a frequency convolution with
the Fej\'er kernel).
\end{theorem}

\begin{theorem}{Aliasing theorem and the Nyquist frequency}{p6-aliasing}
Let $g\in L^1$ be continuous with transform $G$, and let $G$ denote the
transform of the sampled sequence $\{g(t)\}_{t\in\Delta\Ints}$. Under the
summability hypotheses, for every $f$
\[
G(f)=\sum_{k\in\Ints} G(f+k/\Delta).
\]
Sampling \emph{folds} all the continuous-frequency content at
$\{f+k/\Delta\}_{k}$ onto the single frequency $f$. For an SDF this reads
$\sdf(f)=\sum_{k\in\Ints}\sdf(f+k)$ (with $\Delta=1$). The
\textbf{Nyquist frequency} is $1/(2\Delta)$ (i.e.\ $\tfrac12$ when
$\Delta=1$): content above it cannot be recovered.\\[2pt]
\textit{Proof idea:} write $g(t)$ at integer times by both inversion formulas;
split $\int_{-\infty}^{\infty}G(f)e^{2\pi i f t}df$ into period-length blocks,
shift each block by $k/\Delta$ (the phase $e^{2\pi i t k/\Delta}=1$ for
$t\in\Delta\Ints$), swap sum and integral (summability), and match the two
representations of $g(t)$.\\
\textit{Why:} explains why we sample fast enough; if $\sdf$ is negligible
beyond $1/(2\Delta)$ the discrete spectrum faithfully matches the continuous
one, otherwise it is corrupted (poor agreement).
\end{theorem}

\begin{intuition}[Why aliasing matters]
A high-frequency sinusoid sampled too slowly looks identical to a
low-frequency one --- the wagon-wheel effect. Frequencies $f$ and $f+k/\Delta$
are indistinguishable from the samples; they are ``aliases''. The cure is to
sample finely (small $\Delta$, large Nyquist) or to low-pass filter before
sampling.
\end{intuition}

\begin{proposition}{FFT complexity}{p6-fft}
The finite Fourier transform at the $n$ Fourier frequencies $\{k/(\Delta n)\}$
can be computed in $O(n\log n)$ time via the \textbf{Fast Fourier Transform},
versus $O(n^2)$ for the naive sum.\\[2pt]
\textit{Proof idea:} divide-and-conquer --- a transform of length $n$ splits
into two of length $n/2$ (even/odd indices), recursing $\log_2 n$ levels with
$O(n)$ work each.\\
\textit{Why:} makes periodogram/spectral computation feasible at scale; one
can also invert the periodogram by FFT to get the sample ACVS cheaply.
\end{proposition}

\begin{theorem}{Herglotz's theorem}{p6-herglotz}
A complex sequence $\{g_t\}_{t\in\Ints}$ is \textbf{non-negative definite} if
and only if there exists a right-continuous, bounded, non-decreasing function
$G^{(I)}$ with $G^{(I)}(-\tfrac12)=0$ such that
\[
g_t=\int_{-1/2}^{1/2} e^{2\pi i t f}\,dG^{(I)}(f),\qquad t\in\Ints .
\]
\textit{Proof idea:} ($\Leftarrow$) any such integral representation yields a
non-negative-definite sequence because $G^{(I)}$ is non-decreasing.
($\Rightarrow$) construct $G^{(I)}$ as the (weak) limit of the non-negative
Ces\`aro/Fej\'er averages of the partial sums $\sum_{\abs t\le N} g_t
e^{-2\pi i t f}$, which form a tight family of measures.\\
\textit{Why:} it is the existence theorem for the spectrum. Since every ACVS is
non-negative definite (a foundational property of any stationary process),
\emph{every} stationary process has an integrated spectrum $S^{(I)}$ --- even
those, like harmonic processes, with no ordinary SDF.
\end{theorem}

\begin{corollary}{Integrated spectrum exists for every stationary process}{p6-integexists}
Applying Herglotz's theorem to the non-negative-definite ACVS gives
$\acvs_\tau=\int_{-1/2}^{1/2}e^{2\pi i \tau f}\,dS^{(I)}(f)$ with
$S^{(I)}(-\tfrac12)=0$, $S^{(I)}(\tfrac12)=\sigma^2$, and $S^{(I)}$
non-decreasing.\\[2pt]
\textit{Why:} guarantees a frequency-domain description always exists; the SDF
is merely the special (absolutely continuous) case $dS^{(I)}=\sdf\,df$.
\end{corollary}

\begin{theorem}{Lebesgue decomposition of the integrated spectrum}{p6-lebesgue}
Any integrated spectrum decomposes uniquely as
\[
S^{(I)}(\cdot)=S^{(I)}_1(\cdot)+S^{(I)}_2(\cdot)+S^{(I)}_3(\cdot),
\]
where each $S^{(I)}_j$ is non-negative, non-decreasing with
$S^{(I)}_j(-\tfrac12)=0$, and
\begin{enumerate}
  \item $S^{(I)}_1$ is \textbf{absolutely continuous}:
        $S^{(I)}_1(f)=\int_{-1/2}^{f}\sdf(\lambda)\,d\lambda$ (the ordinary SDF
        $\sdf$);
  \item $S^{(I)}_2$ is a \textbf{step function} with jumps $\{p_\ell\}$ at
        frequencies $\{f_\ell\}$ of pure sinusoids (a \emph{line spectrum});
  \item $S^{(I)}_3$ is a \textbf{continuous singular} function (pathological,
        of no practical use).
\end{enumerate}
\textit{Proof idea:} the Lebesgue decomposition theorem for monotone functions
/ measures, specialised to the finite interval $[-\tfrac12,\tfrac12]$.\\
\textit{Why:} classifies the possible spectral behaviours: purely continuous
($S^{(I)}_1$ only), purely discrete/line ($S^{(I)}_2$ only), or mixed.
\end{theorem}

\begin{intuition}[Reading the three components]
\textbf{Absolutely continuous} ($S^{(I)}_1$): the ACVS damps to zero,
$\acvs_\tau\to0$ as $\tau\to\infty$ (since for integrable $\sdf$,
$\int\cos(2\pi f\tau)\sdf\,df\to0$). Example: any ARMA process.
\textbf{Discrete/line} ($S^{(I)}_2$): the ACVS does \emph{not} damp to zero
(a sinusoid keeps oscillating). Example: a sinusoid with random phase and
amplitude (the harmonic process). \textbf{Mixed}: a signal-plus-noise sum of
the two, e.g.\ a periodic component buried in coloured noise.
\end{intuition}

\begin{example}{Spectral characterisation of common processes}{p6-speccases}
\begin{itemize}
  \item \textbf{Purely continuous spectrum} ($S^{(I)}_1\ge0$,
        $S^{(I)}_2\equiv0$): ARMA processes; the SDF $\sdf$ exists and
        $\acvs_\tau\to0$.
  \item \textbf{Purely discrete (line) spectrum} ($S^{(I)}_1\equiv0$,
        $S^{(I)}_2\ge0$): random-phase/amplitude sinusoid; the ACVS persists.
  \item \textbf{Mixed spectrum} ($S^{(I)}_1\ge0$, $S^{(I)}_2\ge0$): pointwise
        aggregation, e.g.\ harmonic signal plus an ARMA background.
\end{itemize}
\end{example}

\begin{theorem}{Spectral representation theorem}{p6-specrep}
Let $\{X_t\}$ be a real-valued discrete-time stationary process with mean
$\mu$. There exists an \textbf{orthogonal-increment process} $\{Z(f)\}$ on
$[-\tfrac12,\tfrac12]$ with
\[
X_t=\mu+\int_{-1/2}^{1/2} e^{2\pi i f t}\,dZ(f)
\quad\text{(equality in mean square),}
\]
where for $f\ne f'$,
\[
\E[dZ(f)]=0,\qquad
\Var\big(dZ(f)\big)=dS^{(I)}(f),\qquad
\Cov\big(dZ(f),dZ(f')\big)=0 .
\]
\textit{Proof idea:} build $Z$ as the $L^2$ (mean-square) limit of finite
Fourier sums of the $X_t$'s; orthogonality of increments mirrors the
non-decrease of $S^{(I)}$, and the variance of an increment is exactly the mass
$S^{(I)}$ places there.\\
\textit{Why:} the deep structural result: \emph{every} stationary process is a
superposition of uncorrelated random sinusoids $e^{2\pi i f t}\,dZ(f)$, with
``amount of variance'' $dS^{(I)}(f)$ at each frequency. It makes precise the
slogan ``$\sdf(f)\,df$ is the power in $[f,f+df]$''.
\end{theorem}

\begin{recallbox}[Recall: covariance of complex random variables]
For complex $Z_1,Z_2$ the covariance is
$\Cov(Z_1,Z_2)=\E[(Z_1-\E Z_1)\conj{(Z_2-\E Z_2)}]$ (note the conjugate on the
second factor). The \emph{complementary covariance} (relation) is
$\mathrm{Rel}\{Z_1,Z_2\}=\E[(Z_1-\E Z_1)(Z_2-\E Z_2)]$ (no conjugate).
\end{recallbox}

\begin{proposition}{Periodogram as squared modulus of $J$}{p6-periogramJ}
With $T=\{1,\dots,n\}$ and $Y_t=X_t-\bar X$,
\[
\hat\sdf^{(p)}(f)=\abs{J(f)}^2,\qquad
J(f)=\sqrt{\tfrac1n}\sum_{t\in T} Y_t\,e^{-2\pi i t f}.
\]
\textit{Proof idea:} starting from
$\hat\sdf^{(p)}(f)=\sum_{\tau}\frac1n\sum_t Y_tY_{t+\abs\tau}e^{-2\pi i\tau f}$,
re-index the double lag sum as a free double sum over $s,t\in T$,
\[
\hat\sdf^{(p)}(f)=\frac1n\sum_{s=1}^{n}\sum_{t=1}^{n}
Y_tY_s\,e^{-2\pi i s f}e^{2\pi i t f}
=\Big(\sqrt{\tfrac1n}\sum_t Y_t e^{-2\pi i t f}\Big)
\conj{\Big(\sqrt{\tfrac1n}\sum_s Y_s e^{-2\pi i s f}\Big)}
=\abs{J(f)}^2 .
\]
\textit{Why:} lets the periodogram be computed by a single FFT in
$O(n\log n)$, and makes manifest that $\hat\sdf^{(p)}\ge0$.
\end{proposition}

\begin{theorem}{Expectation of the periodogram: the Fej\'er kernel}{p6-periogexp}
For a zero-mean process,
\[
\E\big[\hat\sdf^{(p)}(f)\big]
=\sum_{\tau\in\Ints} w_\tau\,\acvs_\tau\,e^{-2\pi i\tau f},\qquad
w_\tau=\Big(1-\frac{\abs\tau}{n}\Big)\ind_{[-n,n]}(\tau),
\]
equivalently the frequency convolution
\[
\E\big[\hat\sdf^{(p)}(f)\big]=\int_{-1/2}^{1/2}\sdf(f')\,\Fej_n(f-f')\,df',
\qquad
\Fej_n(f)=
\begin{cases}
\dfrac1n\dfrac{\sin^2(n\pi f)}{\sin^2(\pi f)}, & f\ne0,\\[2.2ex]
n, & f=0,
\end{cases}
\]
the \textbf{Fej\'er kernel}.\\[2pt]
\textit{Proof idea:} $\E[\hat\acvs_\tau]=\big(1-\tfrac{\abs\tau}{n}\big)\acvs_\tau$
for $\abs\tau<n$; transform term by term. The triangular weight
$w_\tau=\tfrac1n(d*\tilde d)_\tau$ is the autocorrelation of the rectangular
window, whose transform is $\abs{\Dir_T(f)}^2/n=\Fej_n(f)$, so multiplication by
$w_\tau$ in time becomes convolution with $\Fej_n$ in frequency.\\
\textit{Why:} quantifies the \textbf{bias}. Because $\Fej_n$ has wide side
lobes that decay slowly, $\E[\hat\sdf^{(p)}]$ is a \emph{blurred} (leaked)
version of $\sdf$; peaks are smeared and power leaks into neighbouring
frequencies. As $n\to\infty$, $\Fej_n$ concentrates at $0$, so the periodogram
is asymptotically unbiased.
\end{theorem}

\begin{theorem}{Variance of the periodogram: inconsistency}{p6-periogvar}
For $f\ne0\ \mathrm{mod}\ \tfrac12$,
\[
\Var\big(\hat\sdf^{(p)}(f)\big)\xrightarrow[n\to\infty]{}\sdf(f)^2 \;>\;0,
\]
and at distinct Fourier frequencies the ordinates are asymptotically
uncorrelated.\\[2pt]
\textit{Proof idea:} for a Gaussian process $J(f)$ is asymptotically complex
Gaussian, so $\hat\sdf^{(p)}(f)=\abs{J(f)}^2$ is asymptotically
$\tfrac12\sdf(f)\chi^2_2=\sdf(f)\cdot\mathrm{Exp}(1)$, whose variance is
$\sdf(f)^2$.\\
\textit{Why:} the periodogram is \textbf{not consistent} --- its variance does
\emph{not} vanish as $n$ grows; more data merely gives more (equally noisy)
ordinates. This is the motivation for smoothing/averaging (multitapering,
or averaging neighbouring ordinates).
\end{theorem}

\begin{pitfall}
The periodogram is \emph{not} a consistent estimator of the SDF. Doubling $n$
does \textbf{not} halve the variance at a fixed $f$; it stays
$\approx\sdf(f)^2$. ``It got noisier as I added data'' is expected behaviour,
not a bug --- you see finer frequency resolution but the same per-ordinate
scatter. To reduce variance you must \emph{average} (multitaper or smooth).
\end{pitfall}

\begin{theorem}{Expectation of the tapered periodogram}{p6-tapexp}
With a taper $\{h_t\}$ (set $h_t=0$ off $T$) and mean replaced by $\mu$,
\[
\E\big[\hat\sdf^{(p)}_h(f)\big]
=\int_{-1/2}^{1/2}\sdf(f')\,\abs{H(f-f')}^2\,df',
\]
where $H$ is the discrete Fourier transform of $h$.\\[2pt]
\textit{Proof idea:} expand $\hat\sdf^{(p)}_h(f)=\sum_{t,s}h_th_sY_tY_s
e^{-2\pi i f(t-s)}$, take expectations to get $\sum_\tau w_\tau\acvs_\tau
e^{-2\pi i f\tau}$ with $w_\tau=\sum_t h_th_{t+\tau}=(\tilde h*\bar h)_\tau$
(autocorrelation of the taper). By the convolution theorem its transform is
$\tilde H(f)\bar H(f)=\abs{H(f)}^2$, giving the frequency convolution.\\
\textit{Why:} the kernel is now $\abs{H(f)}^2$ instead of the Fej\'er kernel
$\Fej_n$. A well-chosen taper (e.g.\ a dpss / Slepian taper) makes
$\abs{H(f)}^2$ far more concentrated near $0$ with much lower side lobes, so
\textbf{tapering reduces bias / leakage}.
\end{theorem}

\begin{corollary}{Unbiasedness of the tapered periodogram for white noise}{p6-tapunbias}
If $\{X_t\}$ is white noise with variance $\sigma^2$ (so $\sdf\equiv\sigma^2$)
and $\norm{h}_2^2=1$, then $\E[\hat\sdf^{(p)}_h(f)]=\sigma^2=\sdf(f)$ for all
$n$ and all $f$.\\[2pt]
\textit{Proof idea:} the kernel integrates out,
$\E[\hat\sdf^{(p)}_h(f)]=\sigma^2\int_{-1/2}^{1/2}\abs{H(f')}^2\,df'
=\sigma^2\sum_{t}h_t^2=\sigma^2\norm{h}_2^2$ (Parseval), which equals
$\sigma^2$ exactly when $\norm h_2^2=1$.\\
\textit{Why:} pins down the normalisation $\norm h_2^2=1$ in the definition of
a taper, and shows tapering costs nothing in bias for flat spectra.
\end{corollary}

\begin{theorem}{Variance reduction by multitapering}{p6-mtvar}
For $K$ orthonormal tapers, asymptotically
\[
\Cov\big(\hat\sdf^{(p)}_{h_k}(f),\hat\sdf^{(p)}_{h_{k'}}(f)\big)
\to\sdf(f)^2\,\delta_{k,k'},
\quad\text{hence}\quad
\Var\big(\hat\sdf^{(mt)}(f)\big)\approx\frac{\sdf(f)^2}{K}.
\]
\textit{Proof idea:} orthogonality of the tapers makes the $K$ individual
tapered periodograms asymptotically uncorrelated, each with variance
$\sdf(f)^2$; averaging $K$ uncorrelated terms divides the variance by $K$.\\
\textit{Why:} \textbf{multitapering restores consistency}: as $n\to\infty$ one
may let $K\to\infty$ (slowly) so $\Var\to0$. It is the standard fix for the
periodogram's non-vanishing variance, with only a modest bias increase.
\end{theorem}

\begin{intuition}[Bias vs.\ variance, the whole story in one line]
The raw periodogram has tolerable (vanishing) \emph{bias} but stubborn
\emph{variance}. \textbf{Tapering} attacks bias (sharper kernel
$\abs{H}^2$ vs.\ $\Fej_n$); \textbf{multitapering} attacks variance (average
$K$ near-independent estimates, $\Var\approx\sdf^2/K$). Together they give a
usable spectral estimate.
\end{intuition}

% ======================================================================
\subsection{Key Techniques}
% ======================================================================

\begin{keytech}{Compute an SDF from an ACVS}{p6-sdffromacvs}
To get $\sdf(f)=\sum_\tau\acvs_\tau e^{-2\pi i f\tau}$:
\begin{enumerate}
  \item Write down the ACVS (only finitely many lags for MA; geometric for AR).
  \item Pair $\tau$ with $-\tau$ and use
        $e^{-2\pi i f\tau}+e^{2\pi i f\tau}=2\cos(2\pi f\tau)$, giving
        $\sdf(f)=\acvs_0+2\sum_{\tau\ge1}\acvs_\tau\cos(2\pi f\tau)$ (real,
        even).
  \item For a finite linear filter, prefer the factored form: if
        $X_t=\sum_j\psi_j\eps_{t-j}$ with $\eps$ white noise of variance
        $\sigma^2$, then
        $\sdf(f)=\sigma^2\big|\sum_j\psi_j e^{-2\pi i j f}\big|^2$.
\end{enumerate}
For an AR/ARMA $\Phi(B)X_t=\Theta(B)\eps_t$,
$\sdf(f)=\sigma^2\,\dfrac{\abs{\Theta(e^{-2\pi i f})}^2}{\abs{\Phi(e^{-2\pi i f})}^2}$.
\textbf{Sanity checks:} $\sdf$ must be real, even, non-negative, and integrate
to $\acvs_0$ over $[-\tfrac12,\tfrac12]$.
\end{keytech}

\begin{keytech}{Test whether a candidate $\acvs$ is a valid ACVS}{p6-validacvs}
A sequence $\{\acvs_\tau\}$ is a valid ACVS \emph{iff} it is non-negative
definite, which (in the summable case) is equivalent to its Fourier transform
being non-negative:
\[
\sdf(f)=\sum_\tau\acvs_\tau e^{-2\pi i f\tau}\ \ge 0\quad\text{for all }f .
\]
Procedure: (i) check $\acvs_{-\tau}=\acvs_\tau$ and $\acvs_0\ge\abs{\acvs_\tau}$
(necessary); (ii) compute $\sdf(f)$; (iii) hunt for a frequency where
$\sdf(f)<0$ --- often the endpoints $f=0$ (gives $\sum_\tau\acvs_\tau$) or
$f=\tfrac12$ (gives $\sum_\tau(-1)^\tau\acvs_\tau$) are the culprits. A single
negative value \textbf{disqualifies} the candidate.
\end{keytech}

\begin{keytech}{Recover an ACVS from a given SDF}{p6-acvsfromsdf}
Use $\acvs_\tau=\int_{-1/2}^{1/2}\sdf(f)e^{2\pi i f\tau}\,df$. Practical moves:
\begin{itemize}
  \item Drop the imaginary part: since $\acvs_\tau$ is real and $\sdf$ even,
        $\acvs_\tau=\int_{-1/2}^{1/2}\sdf(f)\cos(2\pi f\tau)\,df
        =2\int_0^{1/2}\sdf(f)\cos(2\pi f\tau)\,df$.
  \item For polynomial $\sdf$ integrate by parts; the building blocks are
        $\int_0^{1/2}\cos(2\pi f\tau)\,df=\frac{\sin(\pi\tau)}{2\pi\tau}$ and
        $\int_0^{1/2}f\cos(2\pi f\tau)\,df
        =\frac{\sin(\pi\tau)}{4\pi\tau}+\frac{\cos(\pi\tau)-1}{4\pi^2\tau^2}$.
  \item For a rational $\sdf$ (ARMA), match it to
        $\sigma^2\abs{\Theta}^2/\abs{\Phi}^2$ and read off the recursion, or do
        a partial-fraction / contour integral.
  \item Treat $\tau=0$ separately (the $\sin(\pi\tau)/\tau$ terms need a limit).
\end{itemize}
\end{keytech}

\begin{keytech}{Read off periodogram bias and variance}{p6-readbias}
Given a periodogram problem, recall the three facts:
\begin{itemize}
  \item \textbf{Bias / blurring:} $\E[\hat\sdf^{(p)}]=\sdf*\Fej_n$; the Fej\'er
        kernel's slow-decaying side lobes leak power between frequencies.
        Sharp spectral peaks are the worst affected.
  \item \textbf{Variance:} $\Var(\hat\sdf^{(p)}(f))\to\sdf(f)^2\ne0$ --- not
        consistent.
  \item \textbf{Fixes:} tapering swaps $\Fej_n$ for the sharper $\abs{H}^2$
        (less bias); multitapering averages $K$ tapers
        ($\Var\approx\sdf^2/K$, less variance).
\end{itemize}
\end{keytech}

% ======================================================================
\subsection{Exercises}
% ======================================================================

\begin{exercise}{Fourier transform property table \quad\textnormal{[Sheet 6]}}{p6-ft-props}
For $g\in L^1$ with transform $G$, prove: (1) if $h(t)=\conj{g(t)}$ then
$H(f)=\conj{G(-f)}$; (2) if $h(t)=g(\alpha t)$, $\alpha\ne0$, then
$H(f)=G(f/\alpha)/\abs\alpha$; (3) if $h(t)=g(t+\tau)$ then
$H(f)=G(f)e^{2\pi i f\tau}$; (4) if $h(t)=\alpha g(t)+p(t)$ then
$H(f)=\alpha G(f)+P(f)$.
\end{exercise}
\begin{solution}
\textbf{(1)} Conjugation: pull the conjugate outside (it flips the sign of the
exponent),
\[
H(f)=\int_{-\infty}^{\infty}\conj{g(t)}e^{-2\pi i t f}\,dt
=\conj{\int_{-\infty}^{\infty} g(t)e^{2\pi i t f}\,dt}
=\conj{G(-f)} .
\]
\textbf{(2)} Scaling: substitute $t'=\alpha t$, $dt=dt'/\abs\alpha$ (the
$\abs\alpha$ comes from reversing limits when $\alpha<0$):
\[
H(f)=\int_{-\infty}^{\infty} g(\alpha t)e^{-2\pi i t f}\,dt
=\frac{1}{\abs\alpha}\int_{-\infty}^{\infty} g(t')e^{-2\pi i (f/\alpha)t'}\,dt'
=\frac{1}{\abs\alpha}\,G(f/\alpha).
\]
\textbf{(3)} Shift: insert $1=e^{-2\pi i f\tau}e^{2\pi i f\tau}$,
\[
H(f)=\int_{-\infty}^{\infty} g(t+\tau)e^{-2\pi i f t}\,dt
=\int_{-\infty}^{\infty} g(t+\tau)e^{-2\pi i f(t+\tau)}\,dt\;e^{2\pi i f\tau}
=G(f)e^{2\pi i f\tau}.
\]
\textbf{(4)} Linearity of the integral gives $H=\alpha G+P$ immediately.
\end{solution}

\begin{exercise}{Continuous-time convolution theorem \quad\textnormal{[Sheet 6]}}{p6-convproof}
Let $g,h\in L^1$ and $u=g*h$. Prove $U=G\cdot H$.
\end{exercise}
\begin{solution}
\[
\begin{aligned}
U(f)&=\int_{-\infty}^{\infty} u(x)e^{-2\pi i x f}\,dx
     =\int_{-\infty}^{\infty}\!\!\int_{-\infty}^{\infty} g(x-y)h(y)\,dy\;e^{-2\pi i x f}\,dx\\
    &=\int_{-\infty}^{\infty}\Big(\int_{-\infty}^{\infty} g(x-y)e^{-2\pi i x f}\,dx\Big)h(y)\,dy
     \qquad\text{(Fubini)}\\
    &=\int_{-\infty}^{\infty}\Big(\int_{-\infty}^{\infty} g(s)e^{-2\pi i s f}\,ds\Big)
       e^{-2\pi i y f}h(y)\,dy
     \qquad (s=x-y)\\
    &=G(f)\int_{-\infty}^{\infty} h(y)e^{-2\pi i y f}\,dy=G(f)H(f).
\end{aligned}
\]
The interchange is justified because
$\iint\abs{g(x-y)h(y)}\,dx\,dy=\norm g_1\norm h_1<\infty$, so Fubini applies.
\end{solution}

\begin{exercise}{Aliasing via Poisson summation \quad\textnormal{[Sheet 6]}}{p6-aliasproof}
For continuous $g\in L^1$ with transform $G$, and $G$ the transform of the
sampled sequence $\{g(t)\}_{t\in\Delta\Ints}$, show
$G(f)=\sum_{k\in\Ints} G(f+k/\Delta)$ (under the stated summability
hypotheses).
\end{exercise}
\begin{solution}
For $t\in\Delta\Ints$ both inversion formulas hold:
\[
g(t)=\int_{-\infty}^{\infty} G(f)e^{2\pi i f t}\,df,
\qquad
g(t)=\int_{-1/(2\Delta)}^{1/(2\Delta)} G(f)e^{2\pi i f t}\,df .
\]
Split the real line into period blocks centred at $k/\Delta$ and shift each:
\[
\begin{aligned}
g(t)&=\sum_{k\in\Ints}\int_{k/\Delta-1/(2\Delta)}^{k/\Delta+1/(2\Delta)} G(f)e^{2\pi i f t}\,df
     =\sum_{k\in\Ints}\int_{-1/(2\Delta)}^{1/(2\Delta)} G(f+k/\Delta)e^{2\pi i (f+k/\Delta)t}\,df\\
    &=\sum_{k\in\Ints}\int_{-1/(2\Delta)}^{1/(2\Delta)} G(f+k/\Delta)e^{2\pi i f t}\,df,
\end{aligned}
\]
because $t=z\Delta$ for some $z\in\Ints$ makes $e^{2\pi i (k/\Delta)t}
=e^{2\pi i kz}=1$. By the summability hypothesis
$\sum_k\int\abs{G(f+k/\Delta)}\,df<\infty$ we may swap sum and integral:
\[
\int_{-1/(2\Delta)}^{1/(2\Delta)}\Big(\sum_{k\in\Ints} G(f+k/\Delta)\Big)e^{2\pi i f t}\,df
=g(t)=\int_{-1/(2\Delta)}^{1/(2\Delta)} G(f)e^{2\pi i f t}\,df .
\]
Equality of these Fourier-series coefficients (for all $t\in\Delta\Ints$) forces
$\sum_k G(f+k/\Delta)=G(f)$ for a.e.\ $f$; both sides are continuous (by the
$\ell^1$/summability hypotheses), so equality holds for all $f$.
\end{solution}

\begin{exercise}{SDF of white noise \quad\textnormal{[Sheet 7]}}{p6-wnsdf}
Let $\{\eps_t\}$ be white noise with variance $\sigma_\eps^2$. Find its SDF.
\end{exercise}
\begin{solution}
The ACVS is $\acvs_\tau=\sigma_\eps^2$ for $\tau=0$ and $0$ otherwise, so only
the $\tau=0$ term survives:
\[
\sdf(f)=\sum_{\tau\in\Ints}\acvs_\tau e^{-2\pi i f\tau}=\sigma_\eps^2,
\qquad\text{for all }f .
\]
The spectrum is \textbf{flat} --- equal power at every frequency, which is why
it is called ``white''.
\end{solution}

\begin{exercise}{Is a truncated ACVS a valid ACVS? \quad\textnormal{[Sheet 7]}}{p6-truncacvs}
Let $\acvs_\tau=1$ for $\abs\tau\le K$ and $0$ otherwise, $K$ a positive
integer. Is $\{\acvs_\tau\}$ a valid ACVS? What would the SDF be?
\end{exercise}
\begin{solution}
It is summable and symmetric, so we may form
\[
\sdf(f)=\sum_{\tau=-K}^{K} e^{-2\pi i\tau f}
=1+2\sum_{\tau=1}^{K}\cos(2\pi\tau f)
=\frac{\sin\big((2K+1)\pi f\big)}{\sin(\pi f)}
\]
(the Dirichlet kernel). This is \emph{not} non-negative everywhere. Evaluate at
the Nyquist endpoint $f=\tfrac12$: using $\cos(\pi\tau)=(-1)^\tau$,
\[
\sdf(\tfrac12)=1+2\sum_{\tau=1}^{K}(-1)^\tau=(-1)^K .
\]
So $\sdf(\tfrac12)=-1<0$ whenever $K$ is \textbf{odd} (write $K=2q+1$). Since the
SDF goes negative, $\{\acvs_\tau\}$ is \textbf{not} a valid ACVS for any
stationary process. (It fails to be non-negative definite.)
\end{solution}

\begin{intuition}
Truncating a genuine ACVS to a finite window in general breaks
non-negative-definiteness --- exactly why one cannot just chop the ACVS and
hope. The clean way to ``truncate'' is to taper the spectrum, not the ACVS.
\end{intuition}

\begin{exercise}{ACVS from a triangular SDF \quad\textnormal{[Sheet 7]}}{p6-triangsdf}
Find the ACVS of $\{X_t\}$ with
$\sdf(f)=\dfrac{\sigma^2}{\pi}\,(1-2\abs f)\,\ind_{\{\abs f\le 1/2\}}$.
\end{exercise}
\begin{solution}
For $\tau\ne0$, since $\sdf$ is even the $\sin$ part drops out:
\[
\acvs_\tau=\int_{-1/2}^{1/2} e^{2\pi i f\tau}\sdf(f)\,df
=\frac{2\sigma^2}{\pi}\int_{0}^{1/2}\cos(2\pi f\tau)(1-2f)\,df
=\frac{2\sigma^2}{\pi}\big(I_1-2I_2\big),
\]
with $I_1=\int_0^{1/2}\cos(2\pi f\tau)\,df$ and
$I_2=\int_0^{1/2} f\cos(2\pi f\tau)\,df$. Computing,
\[
I_1=\frac{\sin(\pi\tau)}{2\pi\tau},\qquad
I_2=\frac{\sin(\pi\tau)}{4\pi\tau}+\frac{\cos(\pi\tau)-1}{4\pi^2\tau^2}.
\]
Then $I_1-2I_2=-\dfrac{\cos(\pi\tau)-1}{2\pi^2\tau^2}
=\dfrac{1-\cos(\pi\tau)}{2\pi^2\tau^2}$, so
\[
\boxed{\;\acvs_\tau=\frac{\sigma^2}{\pi^3}\,\frac{1-\cos(\pi\tau)}{\tau^2}\;}
\qquad(\tau\ne0).
\]
For $\tau=0$, $\acvs_0=\dfrac{2\sigma^2}{\pi}\int_0^{1/2}(1-2f)\,df
=\dfrac{2\sigma^2}{\pi}\cdot\tfrac14=\dfrac{\sigma^2}{2\pi}$. (Note
$1-\cos(\pi\tau)=0$ for even $\tau$, $=2$ for odd $\tau$, so the ACVS is
supported on odd lags plus $\tau=0$. Numerically: $\acvs_0=\sigma^2/(2\pi)$,
$\acvs_1=2\sigma^2/\pi^3$, $\acvs_2=0$, $\acvs_3=2\sigma^2/(9\pi^3)$, all
verified.)
\end{solution}

\begin{exercise}{SDF of an MA(q) process \quad\textnormal{[Sheet 7]}}{p6-maqsdf}
What is the SDF of an MA($q$) process
$X_t=\sum_{j=0}^{q}\theta_j\eps_{t-j}$ (here writing $\theta_0=1$ for
concreteness, $\eps$ white noise of variance $\sigma_\eps^2$)?
\end{exercise}
\begin{solution}
Its ACVS is $\acvs_\tau=\sigma_\eps^2\sum_{j=0}^{q}\theta_j\theta_{j+\abs\tau}$
for $\abs\tau\le q$ (zero otherwise), which we write as the double sum
$\acvs_\tau=\sigma_\eps^2\sum_{j=0}^{q}\sum_{k=0}^{q}
\theta_j\theta_k\ind_{\{\tau=k-j\}}$. Since $\acvs\in\ell^1$ the SDF exists:
\[
\begin{aligned}
\sdf(f)&=\sum_{\tau\in\Ints}\acvs_\tau e^{-2\pi i\tau f}
       =\sigma_\eps^2\sum_{j=0}^{q}\sum_{k=0}^{q}\theta_j\theta_k
        \sum_{\tau}\ind_{\{\tau=k-j\}}e^{-2\pi i\tau f}\\
       &=\sigma_\eps^2\sum_{j=0}^{q}\sum_{k=0}^{q}\theta_j\theta_k
        e^{-2\pi i(k-j)f}
       =\sigma_\eps^2\Big(\sum_{j=0}^{q}\theta_j e^{-2\pi i j f}\Big)
        \conj{\Big(\sum_{k=0}^{q}\theta_k e^{-2\pi i k f}\Big)}\\
       &=\sigma_\eps^2\,\Big|\,\sum_{j=0}^{q}\theta_j e^{-2\pi i j f}\,\Big|^2
        =\sigma_\eps^2\,\abs{\Theta(e^{-2\pi i f})}^2 .
\end{aligned}
\]
The SDF is the squared modulus of the MA transfer function --- automatically
non-negative, as any SDF must be.
\end{solution}

\begin{exercise}{Time-varying-amplitude harmonic process \quad\textnormal{[Sheet 7]}}{p6-tvharm}
Let $X_t=\eps_t\cos(\nu t+\Theta)$ with $\nu$ fixed, $\eps_t$ mean-zero white
noise of variance $\sigma^2$ (not constrained positive), independent of the
random phase $\Theta$ whose pdf is
$f_\Theta(\theta)=\tfrac1{2\pi}(1+\cos\theta)$ on $\abs\theta\le\pi$. Find the
first two moments of $X_t$; are they $t$-free?
\end{exercise}
\begin{solution}
\textbf{Mean.} By independence,
$\E[X_t]=\E[\eps_t]\,\E[\cos(\nu t+\Theta)]=0\cdot(\cdot)=0$, constant in $t$
(the second factor is finite since $\abs{\cos}\le1$).\\
\textbf{Autocovariance.} By independence,
$\E[X_tX_{t+\tau}]
=\E[\eps_t\eps_{t+\tau}]\,\E[\cos(\nu t+\Theta)\cos(\nu(t+\tau)+\Theta)]$.
White noise gives $\E[\eps_t\eps_{t+\tau}]=\sigma^2\ind_{\{\tau=0\}}$. Using
$\cos A\cos B=\tfrac12[\cos(A-B)+\cos(A+B)]$,
\[
\E\big[\cos(\nu t+\Theta)\cos(\nu(t+\tau)+\Theta)\big]
=\tfrac12\big[\cos(\nu\tau)+\E[\cos(\nu(2t+\tau)+2\Theta)]\big].
\]
The remaining expectation vanishes for this $\Theta$:
\[
\begin{aligned}
\E[\cos(\nu(2t+\tau)+2\Theta)]
&=\frac1{2\pi}\int_{-\pi}^{\pi}\cos(\nu(2t+\tau)+2\theta)\,d\theta\\
&\quad+\frac1{2\pi}\int_{-\pi}^{\pi}\cos(\nu(2t+\tau)+2\theta)\cos\theta\,d\theta .
\end{aligned}
\]
The first integral is
$\tfrac1{4\pi}\big[\sin(\nu(2t+\tau)+2\theta)\big]_{-\pi}^{\pi}=0$.
For the second, write
$\cos(\cdots+2\theta)\cos\theta=\tfrac12[\cos(\cdots+\theta)+\cos(\cdots+3\theta)]$;
both pieces integrate to $0$ over the full period. Hence
$\E[\cos(\nu(2t+\tau)+2\Theta)]=0$, and
\[
\E[X_tX_{t+\tau}]=\sigma^2\ind_{\{\tau=0\}}\cdot\tfrac12\cos(\nu\tau)
=\begin{cases}\sigma^2/2,&\tau=0,\\[2pt]0,&\tau\ne0.\end{cases}
\]
Both moments are independent of $t$, so $\{X_t\}$ is stationary; in fact it is
white noise with variance $\sigma^2/2$. (The non-uniform phase did not spoil
stationarity precisely because its $2\theta$-harmonics integrate out.)
\end{solution}

\begin{exercise}{Periodogram as $\abs{J(f)}^2$ \quad\textnormal{[Sheet 8]}}{p6-perioJ}
With $T=\{1,\dots,n\}$, show $\hat\sdf^{(p)}(f)=\abs{J(f)}^2$ for
$J(f)=\sqrt{1/n}\sum_{t\in T}(X_t-\bar X)e^{-2\pi i t f}$.
\end{exercise}
\begin{solution}
Put $Y_t=X_t-\bar X$. Starting from the definition and re-indexing the lag sum
as a free double sum over $s,t\in\{1,\dots,n\}$ (the substitution $\tau=t-s$
sweeps all $\abs\tau\le n-1$),
\[
\begin{aligned}
\hat\sdf^{(p)}(f)
&=\sum_{\tau=-n+1}^{n-1}\frac1n\sum_{t=1}^{n-\abs\tau} Y_tY_{t+\abs\tau}e^{-2\pi i\tau f}
 =\frac1n\sum_{s=1}^{n}\sum_{t=1}^{n} Y_tY_s\,e^{-2\pi i s f}e^{2\pi i t f}\\
&=\Big(\sqrt{\tfrac1n}\sum_{t=1}^{n} Y_t e^{-2\pi i t f}\Big)
  \conj{\Big(\sqrt{\tfrac1n}\sum_{s=1}^{n} Y_s e^{-2\pi i s f}\Big)}
 =J(f)\,\conj{J(f)}=\abs{J(f)}^2 .
\end{aligned}
\]
(In the second equality we relabel $s=t+|\tau|$ so the lag $\tau=t-s$ runs
freely, and $e^{-2\pi i\tau f}=e^{2\pi i t f}e^{-2\pi i s f}$ separates the two
conjugate sums.)
\end{solution}

\begin{exercise}{Expectation of the tapered periodogram \quad\textnormal{[Sheet 8]}}{p6-tapexp}
Replacing $\bar X$ by $\mu$, show
$\E[\hat\sdf^{(p)}_h(f)]=\int_{-1/2}^{1/2}\sdf(f')\abs{H(f-f')}^2\,df'$, with
$H$ the transform of $h$.
\end{exercise}
\begin{solution}
With $Y_t=X_t-\mu$ and $g_t=h_t$ ($g_t=0$ off $T$),
\[
\hat\sdf^{(p)}_h(f)=\Big|\sum_{t=1}^{n} g_t Y_t e^{-2\pi i f t}\Big|^2
=\sum_{\tau=-n+1}^{n-1}\sum_{t} g_t g_{t+\abs\tau}\,Y_t Y_{t+\tau}\,e^{-2\pi i f\tau}.
\]
Taking expectations ($\E[Y_tY_{t+\tau}]=\acvs_\tau$),
\[
\E\big[\hat\sdf^{(p)}_h(f)\big]
=\sum_{\tau\in\Ints} w_\tau\,\acvs_\tau\,e^{-2\pi i f\tau},
\qquad w_\tau=\sum_{t\in\Ints} h_t h_{t+\tau}.
\]
By the convolution theorem $\E[\hat\sdf^{(p)}_h(f)]=\int_{-1/2}^{1/2}
\sdf(f')W(f-f')\,df'$ where $W$ is the transform of $w$. Setting
$\tilde h_t=h_{-t}$ and $\bar h_t=h_t$, we recognise $w=\tilde h*\bar h$, so by
the convolution theorem again $W(f)=\tilde H(f)\bar H(f)=\abs{H(f)}^2$.
Therefore
$\E[\hat\sdf^{(p)}_h(f)]=\int_{-1/2}^{1/2}\sdf(f')\abs{H(f-f')}^2\,df'$.
\end{solution}

\begin{exercise}{Unbiased tapered periodogram for white noise \quad\textnormal{[Sheet 8]}}{p6-tunbias}
If $\{X_t\}$ is white noise with variance $\sigma^2$, prove
$\hat\sdf^{(p)}_h(f)$ is unbiased for all $n$ provided $\norm h_2^2=1$.
\end{exercise}
\begin{solution}
White noise has $\sdf\equiv\sigma^2$. Using the tapered-expectation formula and
relabelling (periodicity),
\[
\begin{aligned}
\E\big[\hat\sdf^{(p)}_h(f)\big]
&=\int_{-1/2}^{1/2}\sdf(f-f')\abs{H(f')}^2\,df'
 =\sigma^2\int_{-1/2}^{1/2}\abs{H(f')}^2\,df'\\
&=\sigma^2\sum_{t\in\Ints}\sum_{s\in\Ints} h_t h_s
   \int_{-1/2}^{1/2} e^{-2\pi i(t-s)f'}\,df'
 =\sigma^2\sum_{t} h_t^2=\sigma^2\norm h_2^2,
\end{aligned}
\]
where $\int_{-1/2}^{1/2}e^{-2\pi i\tau f}\,df=\ind_{\{\tau=0\}}$ collapses the
double sum (Parseval). Hence if $\norm h_2^2=1$ then
$\E[\hat\sdf^{(p)}_h(f)]=\sigma^2=\sdf(f)$, i.e.\ exactly unbiased.
\end{solution}

% ---------------------- NEW EXERCISES ----------------------------------

\begin{exercise}{SDF of an MA(1) process \quad\textnormal{[New]}}{p6-ma1sdf}
Let $X_t=\eps_t-\theta\eps_{t-1}$ with $\eps$ white noise of variance
$\sigma^2$. Find $\sdf(f)$ two ways (from the ACVS and from the transfer
function), and locate its maximum and minimum.
\end{exercise}
\begin{solution}
\textbf{From the ACVS.} $\acvs_0=\sigma^2(1+\theta^2)$,
$\acvs_{\pm1}=-\sigma^2\theta$, $0$ otherwise. Then
\[
\sdf(f)=\acvs_0+2\acvs_1\cos(2\pi f)
=\sigma^2\big(1+\theta^2-2\theta\cos(2\pi f)\big).
\]
\textbf{From the transfer function.} $\Theta(z)=1-\theta z$, so
\[
\sdf(f)=\sigma^2\abs{1-\theta e^{-2\pi i f}}^2
=\sigma^2\big(1-\theta e^{-2\pi i f}\big)\big(1-\theta e^{2\pi i f}\big)
=\sigma^2(1+\theta^2-2\theta\cos(2\pi f)),
\]
the same answer. Since $\cos(2\pi f)$ ranges over $[-1,1]$: for $\theta>0$ the
\textbf{minimum} is at $f=0$, $\sdf(0)=\sigma^2(1-\theta)^2$, and the
\textbf{maximum} at $f=\tfrac12$, $\sdf(\tfrac12)=\sigma^2(1+\theta)^2$
(high-pass shape); for $\theta<0$ these swap (low-pass). The SDF is
non-negative everywhere, as required (it touches $0$ only on the unit circle
when $\abs\theta=1$).
\end{solution}

\begin{exercise}{SDF of an AR(1) process \quad\textnormal{[New]}}{p6-ar1sdf}
Let $X_t=\phi X_{t-1}+\eps_t$, $\abs\phi<1$, $\eps$ white noise of variance
$\sigma^2$. Find $\sdf(f)$ from the ACVS and confirm it via the transfer
function. Describe its shape for $\phi>0$ and $\phi<0$.
\end{exercise}
\begin{solution}
The ACVS is $\acvs_\tau=\dfrac{\sigma^2}{1-\phi^2}\,\phi^{\abs\tau}$. Summing the
two geometric tails,
\[
\begin{aligned}
\sdf(f)&=\frac{\sigma^2}{1-\phi^2}\sum_{\tau\in\Ints}\phi^{\abs\tau}e^{-2\pi i f\tau}
        =\frac{\sigma^2}{1-\phi^2}\Big(1+\sum_{\tau\ge1}\phi^\tau
         (e^{-2\pi i f\tau}+e^{2\pi i f\tau})\Big)\\
       &=\frac{\sigma^2}{1-\phi^2}\Big(1+\frac{\phi e^{-2\pi i f}}{1-\phi e^{-2\pi i f}}
         +\frac{\phi e^{2\pi i f}}{1-\phi e^{2\pi i f}}\Big)
        =\frac{\sigma^2}{1-\phi^2}\cdot\frac{1-\phi^2}
         {1-2\phi\cos(2\pi f)+\phi^2}\\
       &=\frac{\sigma^2}{1-2\phi\cos(2\pi f)+\phi^2}.
\end{aligned}
\]
\textbf{Transfer-function check.} $\Phi(z)=1-\phi z$, so
$\sdf(f)=\sigma^2/\abs{1-\phi e^{-2\pi i f}}^2
=\sigma^2/(1-2\phi\cos(2\pi f)+\phi^2)$ --- identical. For $\phi>0$ the
denominator is smallest at $f=0$, so $\sdf$ peaks at low frequency (slowly
varying, ``red'' spectrum); for $\phi<0$ it peaks at $f=\tfrac12$ (rapidly
alternating, ``blue'' spectrum). As $\phi\to1$ the peak at $0$ blows up
(near-unit-root, very long-range dependence).
\end{solution}

\begin{exercise}{Endpoint test for a candidate ACVS \quad\textnormal{[New]}}{p6-endpoint}
Decide whether $\acvs_0=1,\ \acvs_{\pm1}=\rho,\ \acvs_\tau=0\ (\abs\tau\ge2)$ is
a valid ACVS, as a function of $\rho\in\Reals$.
\end{exercise}
\begin{solution}
This is the ACVS shape of an MA(1). Compute the SDF:
\[
\sdf(f)=1+2\rho\cos(2\pi f).
\]
Since $\cos(2\pi f)\in[-1,1]$, the minimum of $\sdf$ is $1-2\abs\rho$, attained
at $f=0$ (if $\rho<0$) or $f=\tfrac12$ (if $\rho>0$). Non-negativity
$\sdf(f)\ge0$ for all $f$ holds \emph{iff} $1-2\abs\rho\ge0$, i.e.
\[
\boxed{\;\abs\rho\le\tfrac12\;}
\]
For $\abs\rho>\tfrac12$ the SDF dips below zero and the sequence is \textbf{not}
a valid ACVS. (Consistent with the fact that an MA(1) cannot have
$\abs{\acf_1}>\tfrac12$.)
\end{solution}

\begin{exercise}{Variance from the SDF; Parseval \quad\textnormal{[New]}}{p6-parseval}
For a stationary process with SDF $\sdf$, show
$\Var(X_t)=\int_{-1/2}^{1/2}\sdf(f)\,df$, and use it to read off $\acvs_0$ for
the AR(1) SDF $\sdf(f)=\sigma^2/(1-2\phi\cos(2\pi f)+\phi^2)$.
\end{exercise}
\begin{solution}
By the inversion formula at lag $\tau=0$,
\[
\acvs_0=\int_{-1/2}^{1/2}\sdf(f)e^{2\pi i f\cdot 0}\,df
=\int_{-1/2}^{1/2}\sdf(f)\,df=\Var(X_t),
\]
the spectral statement of Parseval: total variance equals total spectral mass.
For the AR(1), rather than integrate directly we use the known
$\acvs_0=\sigma^2/(1-\phi^2)$; consistency is the standard integral
\[
\int_{-1/2}^{1/2}\frac{\sigma^2\,df}{1-2\phi\cos(2\pi f)+\phi^2}
=\frac{\sigma^2}{1-\phi^2},
\]
which one can verify by residues (substitute $z=e^{2\pi i f}$; the integrand has
poles at $z=\phi$ and $z=1/\phi$, only $z=\phi$ lying inside the unit circle).
Thus $\Var(X_t)=\sigma^2/(1-\phi^2)$, matching the time-domain computation.
\end{solution}
